% !TEX root = ../main.tex
\begingroup
\linespread{1.20}\selectfont
\begin{center}
  {\yahei\zihao{2} 复杂网络动力学学习与拓扑重构\par}
  \vspace{0.45em}
  {\yahei\zihao{-3} ——从观测时间序列到演化规律与潜在连接\par}
  \vspace{1.4em}
  {\zihao{4} Lrvin-q\par}
  \vspace{0.35em}
  {\zihao{-4} \paperdate\par}
\end{center}

\vspace{0.55em}
\begin{center}
  {\yahei\zihao{3} 摘要}
\end{center}

许多复杂系统都可理解为相互作用的组分在网络上共同演化：节点承载状态，连接传递影响，局部相互作用最终形成可观测的群体行为。
理想的正向模型需要同时知道网络结构、演化规则及其参数，
但真实场景往往只留下各组分的时间序列。
如何从这些观测中学习动力学规律，
并在连接关系不可见时进一步重构潜在拓扑，
是理解、预测和干预复杂系统的基础问题。

本文提出一套由“网络已知”逐步推进到“网络未知”的学习框架。
网络已知时，以分离的自身通道和邻居通道分别表示组分自身演化与交互作用，避免两类信息过早混合；
网络未知时，
在同一动力学学习器外引入离散拓扑生成器，
通过硬图上的规律学习、软图上的结构更新和结构信息更充分的样本组织，
减弱错误连接与动力学参数相互补偿的现象。
为使规律学习、参数诊断和拓扑恢复都能被清楚检验，
本文选取复杂网络 SIR 传播过程作为代表性案例。
该模型既可描述传染病传播，
也可作为谣言、信息等扩散现象的简化原型。
实验在 ER、WS 和 BA 三类网络上展开。

\begin{itemize}[leftmargin=2.3em,itemsep=0.05em,topsep=0.2em]
  \item {\yahei 网络已知时，分离自身演化与邻居交互的双通道模型在三类网络上取得 0.0130 的平均感染概率曲线误差、0.0041 的感染率反演绝对误差和 0.0079 的感染者占比自回归误差，
说明合适的交互表示有助于同时保持预测精度与规律可解释性。}
  \item {\yahei 网络未知时，所选模型的硬图 F1 达到 0.937--1.000，感染率反演值为 0.289--0.300，
感染者占比自回归误差不超过 0.0083，
表明仅凭状态时间序列也有可能同时恢复有效的演化规律与主要连接。}
  \item {\yahei 联合学习在拓扑恢复、参数反演与多步预测三个层面取得一致结果；
最低验证交叉熵在本组多初始化实验中选出高质量模型，
表明结构学习与动力学学习能够相互支持。}
\end{itemize}

本文以 SIR 案例验证了把网络动力学分解为自身演化与交互作用、并联合学习结构和规律的研究路径。
该框架为其他传播模型与连续网络动力系统的数据驱动建模提供了可进一步检验的方法基础。

\noindent{\yahei 关键词：}复杂系统；
复杂网络；
动力学学习；
拓扑重构；
时间序列；
图神经网络；
SIR 传播
\endgroup
